\chapter*{卷四\quad 后题}
\phantomsection\addcontentsline{toc}{chapter}{卷四\quad 后题}
\markboth{卷四\quad 后题}{卷四\quad 后题}

卷四六片，均非城池尚在运转之时所刻，或为后世题记，或为伪托，或为抄手之讹。此卷之所以收入本书，唯因其同在原件之中；况且一片文字若非当时所刻，其史料价值反而倍增，盖其证明此后仍有后来者曾抵此地。

\begin{jian}{简三十二}
后三百七十二年，有客至。门不可启，壁上题名而去。

题名者三人，其一名不可辨。
\end{jian}
\chuchu{原件残片 81 号，长 13.4 cm，金属，字口浅，与他片不类。}
\xiaojì{「后三百七十二年」若自前 221 年起算，则为公元 151 年，其时中原当东汉桓帝年间。「有客至」不称职衔，亦不留姓字，唯记「题名而去」，所题三人之名而一不可辨，此三人为谁全书无据。「门不可启」四字与简二十一「第十九项：门，闭」遥相呼应，足见历三百七十二年，其门犹然未启。}

\begin{jian}{简三十三}
此非月也，此舟也。舟已不行。吾辈守之三世矣。

守者无禄，亦无代。
\end{jian}
\chuchu{原件残片 83 号，长 18.1 cm，金属。}
\xiaojì{「此舟也」一语，与笔者《Лунная прародина: опыт реконструкции》之基本判断相合；然此片字迹拙涩，刻痕深浅不均，疑为后世摹刻而非原件，此类残片之性质本书无从判定，详下编第七章末。「吾辈守之三世矣」，则守城者非一人一世而至少历三代；历三代守一座已闭之城而无禄无代，其缘由不可深考。「无代」二字，或谓无人接替，或谓不以他物相代，二义俱不可定。}

\begin{jian}{简三十四}
大明万历十七年，内官监奉旨采办，至此而返。

同行者二百七十人，还者一百四十一人。
\end{jian}
\chuchu{原件残片 85 号，长 11.2 cm，陶瓷，背面有朱印一方，不可辨。}
\xiaojì{「大明万历十七年」即一五八九年，乃全书最晚之一条自记年代。「至此而返」之「此」字何指无从判定，若指广寒宫，则明代尚有人登临月面之说始见于此；然「内官监」「采办」皆明代宫廷惯用语汇，见于此处颇觉辞气不谐，疑出后人附会。至「同行者二百七十人，还者一百四十一人」，近半数不还，此数与万历朝任何已知远航记录皆不吻合，无从核校。}

\begin{jian}{简三十五}
□有十二，不可尽。尽者三，余九亡。

十二者，三之四倍。三尽而九亡，亡者不留。
\end{jian}
\chuchu{原件残片 87 号，长 9.8 cm，金属，上下俱蚀。}
\xiaojì{「□有十二」，首字蚀尽而无据可补。「十二」与玉山十二个乘员舱之数相合，然通篇并无一处可与乘员舱明确对应；「尽者三，余九亡」，三与九之和恰为十二，此数尤须留意。第二行「十二者，三之四倍」以下乃他片绝无之推演算语，与全书其余各片之记账笔法迥异，疑系后世注文误入正文。}

\begin{jian}{简三十六}
一物，圆，径不及二寸，重若干。非金非石。击之不鸣。来源不明，暂不归类。
\end{jian}
\chuchu{原件残片 89 号，长 12.6 cm，金属。此片所述之圆片，现存基金会，编号 ФВ-1963-46。}
\xiaojì{「非金非石」「击之不鸣」，其性状描写与基金会 \textsc{фв}-1963-46 编目卷宗所收金属圆片相合；该器物经俄方检测，铼锇同位素比值异常，地球无对照矿源。「暂不归类」四字实为一九六三年编目时之卡片批注，非原件刻文，乃写入编目卡而非刻于残片，此为全书编次之一处特例。「重若干」三字原件如此，似为尚未填实之公文格式。}

\begin{jian}{简三十七}
字尽于此。

此后无字。
\end{jian}
\chuchu{原件残片 91 号，长 6.2 cm，金属，为一残片之末段，前文俱失。}
\xiaojì{「字尽于此」既可为抄手自记，亦可为全书之殿尾，以本书编次论取后一义。「此后无字」四字系后人补刻，字口与上文有别，与他片所见补刻同出一手；补刻此四字者，显系自知此片乃全卷之末片。}
